\documentclass[11pt]{article}

\usepackage[margin=1in]{geometry}
\usepackage{amsmath,amssymb,bm}
\usepackage{hyperref}

\title{Theory Behind \texttt{GPUSolver::smallestCondensedMechanicalEigenpair}}
\author{}
\date{}

\begin{document}

\maketitle

\section{Purpose}

The function \texttt{GPUSolver::smallestCondensedMechanicalEigenpair}
computes the smallest eigenvalue and corresponding eigenvector of the
mechanical stability problem after eliminating the electric potential
degrees of freedom. In the flexoelectric problem, the unknown vector is
ordered as
\[
    \bm z =
    \begin{bmatrix}
        \bm u \\
        \bm \phi
    \end{bmatrix},
\]
where \(\bm u\) contains the displacement degrees of freedom and
\(\bm \phi\) contains the electric potential degrees of freedom.

At a converged nonlinear state, the assembled tangent/Hessian matrix has
the block form
\[
    \bm H =
    \begin{bmatrix}
        \bm H_{uu} & \bm H_{u\phi} \\
        \bm H_{\phi u} & \bm H_{\phi\phi}
    \end{bmatrix}.
\]

The goal is to test mechanical stability while allowing the electric
potential field to relax consistently with the mechanical perturbation.

\section{Condensation of the Electric Potential}

For a perturbation \((\delta \bm u, \delta \bm \phi)\), the quadratic
second variation is
\[
    \delta^2 \Pi =
    \begin{bmatrix}
        \delta \bm u \\
        \delta \bm \phi
    \end{bmatrix}^{T}
    \begin{bmatrix}
        \bm H_{uu} & \bm H_{u\phi} \\
        \bm H_{\phi u} & \bm H_{\phi\phi}
    \end{bmatrix}
    \begin{bmatrix}
        \delta \bm u \\
        \delta \bm \phi
    \end{bmatrix}.
\]

For a given mechanical perturbation \(\delta \bm u\), the compatible
electric perturbation is obtained by enforcing stationarity with respect
to \(\delta \bm \phi\):
\[
    \bm H_{\phi u}\delta \bm u
    +
    \bm H_{\phi\phi}\delta \bm \phi
    =
    \bm 0.
\]
Assuming \(\bm H_{\phi\phi}\) is nonsingular,
\[
    \delta \bm \phi
    =
    -\bm H_{\phi\phi}^{-1}\bm H_{\phi u}\delta \bm u.
\]

Substituting this expression into the quadratic form gives the condensed
mechanical Hessian
\[
    \widehat{\bm H}
    =
    \bm H_{uu}
    -
    \bm H_{u\phi}
    \bm H_{\phi\phi}^{-1}
    \bm H_{\phi u}.
\]

Therefore the condensed mechanical eigenvalue problem is
\[
    \widehat{\bm H}\bm x = \lambda \bm x.
\]
The smallest eigenvalue \(\lambda_{\min}\) is used as the stability
indicator. When it crosses zero, the current equilibrium loses stability
in the condensed mechanical sense.

\section{Computing the Smallest Eigenvector by Inverse Iteration}

The purpose of this section is to explain how the code updates the
mechanical eigenvector during the outer iteration. At this stage, assume
the code already has a normalized mechanical vector
\[
    \bm x_k.
\]
The goal of one inverse-iteration step is to compute
\[
    \bm y_k = \widehat{\bm H}^{-1}\bm x_k.
\]
If this operation can be performed repeatedly, then the direction of
\(\bm x_k\) will converge to the eigenvector associated with the
eigenvalue of \(\widehat{\bm H}\) closest to zero.

\subsection*{Applying \(\widehat{\bm H}^{-1}\) Without Forming It}

The code does not explicitly form \(\widehat{\bm H}^{-1}\). It also does
not need to explicitly form a dense inverse of the Schur complement.
Instead, it applies \(\widehat{\bm H}^{-1}\) to the current vector
\(\bm x_k\) by solving one full coupled linear system.

The right-hand side of this coupled solve is
\[
    \begin{bmatrix}
        \bm x_k \\
        \bm 0
    \end{bmatrix}.
\]
The upper part is the current mechanical vector, and the lower part is
zero because we are applying a purely mechanical load in the condensed
mechanical problem.

The code solves
\[
    \begin{bmatrix}
        \bm H_{uu} & \bm H_{u\phi} \\
        \bm H_{\phi u} & \bm H_{\phi\phi}
    \end{bmatrix}
    \begin{bmatrix}
        \bm y \\
        \bm \eta
    \end{bmatrix}
    =
    \begin{bmatrix}
        \bm x_k \\
        \bm 0
    \end{bmatrix}.
\]
The unknown \(\bm y\) is the displacement part of the coupled solution.
The unknown \(\bm \eta\) is the electric-potential part needed to satisfy
the coupled system during this solve.

The block equations are
\[
    \bm H_{uu}\bm y + \bm H_{u\phi}\bm \eta = \bm x_k,
\]
\[
    \bm H_{\phi u}\bm y + \bm H_{\phi\phi}\bm \eta = \bm 0.
\]
The second equation can be solved for the electric part:
\[
    \bm \eta =
    -\bm H_{\phi\phi}^{-1}\bm H_{\phi u}\bm y.
\]
Substituting this expression into the first equation removes
\(\bm \eta\):
\[
    \left(
        \bm H_{uu}
        -
        \bm H_{u\phi}
        \bm H_{\phi\phi}^{-1}
        \bm H_{\phi u}
    \right)\bm y
    =
    \bm x_k.
\]
The matrix in parentheses is exactly the condensed mechanical Hessian
\(\widehat{\bm H}\). Therefore
\[
    \widehat{\bm H}\bm y = \bm x_k.
\]
Thus
\[
    \bm y = \widehat{\bm H}^{-1}\bm x_k.
\]
This shows why solving the full coupled system gives the same
mechanical vector that would be obtained by applying the inverse
condensed mechanical operator.

\subsection*{Updating the Eigenvector Direction}

After the coupled solve, the code keeps only the displacement part
\(\bm y\). This vector is the inverse-iteration update direction. The
code then normalizes it:
\[
    \bm x_{k+1}
    =
    \frac{\bm y}{\|\bm y\|}.
\]
Therefore one outer iteration is
\[
    \bm y_k = \widehat{\bm H}^{-1}\bm x_k,
    \qquad
    \bm x_{k+1} = \frac{\bm y_k}{\|\bm y_k\|}.
\]
This is inverse iteration applied to the condensed mechanical Hessian.

\subsection*{Why This Finds the Eigenvector Closest to Zero}

If \(\widehat{\bm H}\) has eigenpairs
\[
    \widehat{\bm H}\bm q_i = \lambda_i \bm q_i,
\]
then
\[
    \widehat{\bm H}^{-1}\bm q_i =
    \frac{1}{\lambda_i}\bm q_i.
\]
To see the amplification more explicitly, write the current iterate as a
linear combination of eigenvectors:
\[
    \bm x_k = \sum_i c_i^{(k)} \bm q_i .
\]
Here \(c_i^{(k)}\) measures how much of eigenvector \(\bm q_i\) is
present in the current iterate. Applying the inverse operator gives
\[
    \bm y_k
    =
    \widehat{\bm H}^{-1}\bm x_k
    =
    \sum_i c_i^{(k)}
    \widehat{\bm H}^{-1}\bm q_i
    =
    \sum_i
    \frac{c_i^{(k)}}{\lambda_i}
    \bm q_i .
\]
Thus each eigenvector component is multiplied by the factor
\[
    \frac{1}{\lambda_i}.
\]
This is the central reason inverse iteration works. Components
corresponding to smaller \(|\lambda_i|\) are multiplied by larger
\(|1/\lambda_i|\). Suppose \(\lambda_1\) is the eigenvalue closest to
zero, so that
\[
    |\lambda_1| < |\lambda_j|, \qquad j \ne 1.
\]
Then after one inverse application, the relative size of the
\(\bm q_1\) component compared with the \(\bm q_j\) component changes as
\[
    \left|
    \frac{c_1^{(k)}/\lambda_1}
         {c_j^{(k)}/\lambda_j}
    \right|
    =
    \left|
    \frac{c_1^{(k)}}{c_j^{(k)}}
    \right|
    \left|
    \frac{\lambda_j}{\lambda_1}
    \right|.
\]
Because \(|\lambda_j/\lambda_1|>1\), the relative contribution of
\(\bm q_1\) grows after each inverse iteration. In other words, even if
\(\bm q_1\) is only a small part of the initial vector, it is magnified
faster than the other eigenvector components. After \(m\) ideal inverse
applications,
\[
    \bm x_{k+m}
    \quad \hbox{is dominated by} \quad
    \bm q_1,
\]
provided the initial vector has a nonzero component in the
\(\bm q_1\) direction.

This gives the main requirement on the initial vector. If
\[
    \bm x_0 = \sum_i c_i^{(0)}\bm q_i,
\]
then convergence to \(\bm q_1\) requires
\[
    c_1^{(0)} \ne 0.
\]
If \(c_1^{(0)}=0\) exactly, inverse iteration cannot create that missing
component, because each application of \(\widehat{\bm H}^{-1}\) only
rescales the coefficients already present:
\[
    \widehat{\bm H}^{-1}\bm x_0
    =
    \sum_i
    \frac{c_i^{(0)}}{\lambda_i}\bm q_i.
\]
Therefore the \(\bm q_1\) coefficient remains zero for all later
iterations. In practice, this exact orthogonality is unlikely if a
generic or random-looking initial vector is used. The code initializes
\(\bm x_0\) with a deterministic sine--cosine pattern, which is intended
to give a nonzero projection onto the critical eigenvector while avoiding
the need for random number generation.

The normalization step,
\[
    \bm x_{k+1} = \frac{\bm y_k}{\|\bm y_k\|}
\]
does not change which eigenvector components are relatively larger or
smaller. It only prevents the vector norm from growing without bound.
Therefore the direction of the iterate converges to the eigenvector
whose eigenvalue has the largest reciprocal magnitude, which is the
eigenvalue closest to zero in magnitude.

In a stability calculation, this mode is important because loss of
stability occurs when a mechanical eigenvalue approaches and crosses
zero. The inverse iteration therefore naturally focuses on the critical
mode near loss of stability.

\section{Eigenvalue Estimate During Iteration}

During inverse iteration, the code needs a scalar estimate of the
eigenvalue so it can decide whether the outer iteration has converged.
At each iteration, the code already has two vectors:
\[
    \bm x,
\]
the current normalized displacement vector, and
\[
    \bm y,
\]
the result of applying the inverse condensed operator to \(\bm x\).
That is,
\[
    \bm y = \widehat{\bm H}^{-1}\bm x,
\]
where \(\bm y\) is obtained from the full coupled solve described in
Section 3.

The code then computes
\[
    \mu = \bm x^T\bm y
\]
and stores this value in the variable \(\texttt{inverseEigenvalue}\):
\[
    \texttt{inverseEigenvalue} = \texttt{x.dot(y)}.
\]
This quantity is called \(\texttt{inverseEigenvalue}\) because it is an
estimate of \(1/\lambda\), not \(\lambda\) itself.

To see why, first consider the ideal case where the current vector is
exactly an eigenvector:
\[
    \bm x = \bm q,
    \qquad
    \widehat{\bm H}\bm q = \lambda \bm q,
    \qquad
    \|\bm q\|=1.
\]
Applying the inverse condensed operator gives
\[
    \bm y =
    \widehat{\bm H}^{-1}\bm q =
    \frac{1}{\lambda}\bm q.
\]
Then the dot product is
\[
    \mu
    =
    \bm q^T
    \left(
        \frac{1}{\lambda}\bm q
    \right)
    =
    \frac{1}{\lambda},
\]
because \(\bm q^T\bm q=1\). Therefore
\[
    \lambda = \frac{1}{\mu}.
\]

In the actual iteration, \(\bm x\) is not exactly equal to the
eigenvector until convergence. However, after several inverse-iteration
steps, \(\bm x\) becomes close to the critical eigenvector. Then
\[
    \bm y
    =
    \widehat{\bm H}^{-1}\bm x
    \approx
    \frac{1}{\lambda}\bm x,
\]
so
\[
    \bm x^T\bm y
    \approx
    \frac{1}{\lambda}\bm x^T\bm x.
\]
Since the code normalizes \(\bm x\), we have
\[
    \bm x^T\bm x = 1,
\]
and therefore
\[
    \bm x^T\bm y \approx \frac{1}{\lambda}.
\]
This gives the iteration-level eigenvalue estimate
\[
    \lambda \approx \frac{1}{\bm x^T\bm y}.
\]
In the code, this appears as
\[
    \texttt{smallestEigenvalue}
    =
    \frac{1}{\texttt{inverseEigenvalue}}.
\]

This estimate is used only to monitor convergence during inverse
iteration. The code compares the current estimate with the previous one:
\[
    |\lambda_k-\lambda_{k-1}|.
\]
The iteration stops when this change is small relative to the current
eigenvalue scale:
\[
    |\lambda_k-\lambda_{k-1}|
    <
    \texttt{tolerance}
    \max(1,|\lambda_k|).
\]
After the eigenvector has converged, the code recomputes the eigenvalue
using the Rayleigh quotient described in the next section.

\section{Final Rayleigh Quotient}

After convergence, the function recomputes the eigenvalue using the
Rayleigh quotient of the condensed Hessian:
\[
    \lambda
    =
    \bm x^T \widehat{\bm H}\bm x.
\]
The code evaluates this without explicitly storing the dense Schur
complement:
\[
    \widehat{\bm H}\bm x
    =
    \bm H_{uu}\bm x
    -
    \bm H_{u\phi}
    \bm H_{\phi\phi}^{-1}
    \bm H_{\phi u}\bm x.
\]

The important point is that the Rayleigh quotient only needs the vector
\(\widehat{\bm H}\bm x\). It does not need every entry of
\(\widehat{\bm H}\). Therefore the code evaluates the product
\(\widehat{\bm H}\bm x\) from right to left.

First, it computes the electric residual produced by the mechanical
eigenvector:
\[
    \bm r = \bm H_{\phi u}\bm x.
\]
In the code, this is
\[
    \texttt{hpuX} = \bm H_{\phi u}\bm x.
\]

Second, it solves the electric block system
\[
    \bm H_{\phi\phi}\bm w = \bm r.
\]
Thus
\[
    \bm w
    =
    \bm H_{\phi\phi}^{-1}\bm H_{\phi u}\bm x.
\]
In the code, this is
\[
    \texttt{invHppHpuX}
    =
    \bm H_{\phi\phi}^{-1}\bm H_{\phi u}\bm x.
\]

Third, it forms the condensed mechanical matrix-vector product
\[
    \bm z
    =
    \bm H_{uu}\bm x
    -
    \bm H_{u\phi}\bm w.
\]
Since \(\bm w=\bm H_{\phi\phi}^{-1}\bm H_{\phi u}\bm x\), this is exactly
\[
    \bm z
    =
    \widehat{\bm H}\bm x.
\]
In the code, this vector is named \(\texttt{hhatX}\):
\[
    \texttt{hhatX}
    =
    \bm H_{uu}\bm x
    -
    \bm H_{u\phi}\texttt{invHppHpuX}.
\]

Finally, the Rayleigh quotient is computed as
\[
    \lambda
    =
    \bm x^T\bm z
    =
    \bm x^T\widehat{\bm H}\bm x.
\]
In the code, this is
\[
    \texttt{smallestEigenvalue}
    =
    \texttt{x.dot(hhatX)}.
\]

This treatment is different from the inverse iteration step in Section
3. In Section 3, the code applies the inverse condensed operator,
\[
    \widehat{\bm H}^{-1}\bm x,
\]
by solving the full coupled system with right-hand side
\([\,\bm x;\bm 0\,]\). That solve produces the next iterate direction
\(\bm y=\widehat{\bm H}^{-1}\bm x\). In the final Rayleigh quotient, the
code applies the condensed operator itself,
\[
    \widehat{\bm H}\bm x,
\]
using two block matrix-vector products and one solve with the electric
block \(\bm H_{\phi\phi}\). Therefore Section 3 uses the full coupled
solve to move the eigenvector iteration forward, while Section 5 uses
the Schur-complement formula to evaluate the eigenvalue of the already
converged vector.

This final Rayleigh quotient is usually more directly tied to the
condensed mechanical operator than the reciprocal estimate used during
the iteration.

\section{Recovered Full Eigenvector}

The displacement eigenvector returned by the function is
\[
    \delta \bm u = \bm x.
\]
The associated electric potential eigenvector is recovered from the
condensation relation:
\[
    \delta \bm \phi
    =
    -\bm H_{\phi\phi}^{-1}
    \bm H_{\phi u}\bm x.
\]
Therefore the full coupled perturbation is
\[
    \delta \bm z =
    \begin{bmatrix}
        \delta \bm u \\
        \delta \bm \phi
    \end{bmatrix}
    =
    \begin{bmatrix}
        \bm x \\
        -\bm H_{\phi\phi}^{-1}\bm H_{\phi u}\bm x
    \end{bmatrix}.
\]

\section{Assumptions}

The procedure relies on the following assumptions:
\begin{itemize}
    \item The global tangent matrix is ordered as displacement blocks
          followed by the electric potential block.
    \item The electric block \(\bm H_{\phi\phi}\) is nonsingular.
    \item The condensed mechanical Hessian
          \(\widehat{\bm H}\) is the correct stability operator for the
          chosen electrical boundary conditions.
    \item The eigenvalue of interest is the one closest to zero, so that
          inverse iteration on \(\widehat{\bm H}^{-1}\) converges to the
          critical mode.
    \item The initial vector is nonzero, finite, and has a nonzero
          projection onto the desired critical eigenvector. If the
          critical eigenvalue has multiplicity greater than one, the
          method may converge to a linear combination inside the
          corresponding eigenspace.
\end{itemize}

\section{Connection to the Implementation}

The implementation follows these mathematical steps:
\begin{enumerate}
    \item Assemble the full tangent matrix \(\bm H\).
    \item Split it into
          \(\bm H_{uu}\), \(\bm H_{u\phi}\),
          \(\bm H_{\phi u}\), and \(\bm H_{\phi\phi}\).
    \item Factorize the full coupled matrix for inverse iteration.
    \item Repeatedly solve
          \[
              \bm H
              \begin{bmatrix}
                  \bm y \\
                  \bm \eta
              \end{bmatrix}
              =
              \begin{bmatrix}
                  \bm x \\
                  \bm 0
              \end{bmatrix}.
          \]
    \item Keep only the displacement part \(\bm y\), normalize it, and
          update \(\bm x\).
    \item Estimate the eigenvalue during iteration as
          \[
              \lambda \approx \frac{1}{\bm x^T\bm y}.
          \]
    \item After convergence, compute the final value using
          \[
              \lambda = \bm x^T\widehat{\bm H}\bm x.
          \]
    \item Return \(\bm x\) as the displacement eigenvector and recover
          the electric part using
          \[
              \delta\bm\phi
              =
              -\bm H_{\phi\phi}^{-1}\bm H_{\phi u}\bm x.
          \]
\end{enumerate}

\end{document}
